\chapter{Audit of the Physical Claims}
\label{ch:physical-claims}

\section{Planck force and force unification}

The quantity

\begin{equation}
  F_{\mathrm{P}} = \frac{c^4}{G}
\end{equation}

has dimensions of force. It can also be written as the Planck energy divided by
the Planck length. These are algebraic identities among constants and derived
units \citep{planck1899natural}.

\begin{evidence}{Source claim}
The reviewed Superforce manuscript interprets this scale as a force of
unification \citep{pais2023superforce}. The repository reproduces the
dimensional cancellation and the value of the scale.
\end{evidence}

\begin{evidence}{Assessment: unsupported extrapolation}
Dimensional agreement is necessary for an equation but does not derive an
interaction. A force-unification claim requires a field content, symmetry,
couplings, renormalization behavior, and predictions that differ from general
relativity and the Standard Model. None follows from $c^4/G$ alone. Substituting
Planck units into Newtonian or Coulomb force laws selects a scale; it does not
show that distinct gauge interactions become the same mechanism.
\end{evidence}

The absence of $\hbar$ from $c^4/G$ also has no independent physical content:
$\hbar$ cancels in the ratio of two Planck units. The cancellation cannot, by
itself, establish a macroscopic quantum effect.

\section{Root-lattice charges and fluid triads}

The earlier manuscript proposed that Fourier modes inherit a class in the
discriminant group $P/Q$ and that triad closure conserves those classes. The
proposal fails at its first stated map. If the forcing vectors are roots, then
they lie in the root lattice $Q$ and every root represents the zero class in
$P/Q$. The proposed nontrivial parity charge therefore does not arise from the
root set as written.

A nontrivial construction would need, at minimum:

\begin{enumerate}
  \item a homomorphism from the physical Fourier lattice to $P/Q$;
  \item a reason physical modes occupy nonzero weight-lattice cosets;
  \item a forcing or interaction operator that enforces the proposed charge;
  \item a count of admitted and rejected triads under that operator; and
  \item a control forcing with matched spectrum, power, and anisotropy but no
    quotient structure.
\end{enumerate}

\begin{evidence}{Assessment: mechanism not implemented}
The repository contains prose describing a modular filter, but the simulation
step contains no $P/Q$ labels and no triad-admission test. The discriminant-group
argument cannot explain a numerical result until that missing map and operator
exist.
\end{evidence}

\section{The beta-plane interpretation}

The Rhines scale follows from competition between turbulent advection and
Rossby-wave dynamics on a beta plane \citep{rhines1975waves,vallis1993generation}.
A beta-plane simulation therefore needs a beta-dependent term in its dynamical
equations, not merely a spectrum that can be assigned a preferred wavenumber
after the run.

The checked-in LBM parameter object has no beta parameter. Its CPU step performs
equilibrium evaluation, BGK collision, streaming, and macroscopic update. The
accelerated step adds a spatially varying relaxation time and random noise, but
no beta-plane body force or potential-vorticity gradient. The high-resolution
driver sets grid size, time steps, relaxation time, and harmonic amplitude.
It does not configure beta.

The implementation selects the first $N$ roots and weights their absolute first
two coordinates by enumeration index. Root generation now uses a deterministic
lexicographic order, removing a reproducibility defect. Determinism does not
make the construction Weyl invariant. The retained order audit finds that
reversal is exactly invariant because lexicographic reversal pairs opposite
roots with equal absolute coordinates, while a one-place cyclic shift changes
the initialized perturbation by 0.131 in relative $L^2$. Any observed
anisotropy remains confounded with this ordering choice, projection, and index
weighting until a symmetry-invariant map and matched controls remove those
alternatives.

The earlier diagnostic

\begin{equation}
  \beta_{\mathrm{inferred}} = U_{\star} k_{\mathrm{jet}}^2
\end{equation}

is useful only after a beta-plane model and an independently defined velocity
scale exist. Choosing or reconstructing $U_{\star}$ from the same output and
then declaring agreement with an asserted beta is circular.

The retained final high-resolution snapshot also contradicts the published
numbers in the superseded assembly. A direct read of the step-500 snapshot gives
a mean speed near $6.92\times10^{-3}$, a maximum speed near
$1.10\times10^{-2}$, and dominant zonal-mean Fourier mode 17. The zonal-mean
RMS is only $8.75\times10^{-6}$, or $1.22\times10^{-3}$ of total velocity RMS,
and the vorticity spectral second-moment ratio is 0.992. Thus the selected peak
mode is not evidence of a strong zonal state. The earlier text used
$U\approx0.6$ and mode 4. The run driver does not compute a jet count, spectral
slope, or uncertainty interval, so those earlier values have no traceable
analysis path.

\begin{table}[htbp]
  \centering
  \caption{Requirements for a defensible beta-plane comparison.}
  \label{tab:beta-requirements}
  \begin{tabularx}{\textwidth}{>{\raggedright\arraybackslash\bfseries}p{0.28\textwidth}YY}
    \toprule
    Requirement & Repository state & Consequence \\
    \midrule
    Beta-dependent dynamics & Absent from the retained LBM steps & The run is not a beta-plane experiment. \\
    Independent input beta & No beta field in the parameter object & Input-output agreement cannot be tested. \\
    Prespecified jet diagnostic & Not retained with uncertainty or mode-selection rule & Jet count is post hoc. \\
    Matched controls & No isotropic or spectrum-matched ablation & Algebraic causation is unidentified. \\
    Conservation check & JAX drift is $5.60\times10^{-6}$ at step 500; the repaired CPU regression is conservative & Catastrophic growth does not describe the high-resolution run. \\
    \bottomrule
  \end{tabularx}
\end{table}

\section{Affine algebras and conformal invariance}

For an affine Lie algebra at level $k$, the Sugawara construction produces a
Virasoro action under specified representation-theoretic conditions. For
affine $E_8$ at level one, the central charge is

\begin{equation}
  c = \frac{k\,\dim E_8}{k+h^{\vee}} = \frac{248}{31}=8.
\end{equation}

This is a statement about a current algebra and its conformal field theory. It
does not imply that a fluid forced with a pattern labelled $E_8$ inherits that
Virasoro representation. Two-dimensional turbulence can display conformal
statistics in appropriate regimes \citep{bernard2006conformal}, but similarity
of vocabulary is not a derivation.

\begin{evidence}{Assessment: category error}
Persistent homology can summarize geometry in data. It cannot prove that a
fluid field carries a Sugawara stress tensor. The latter requires an explicit
current algebra, operator products or equivalent observables, and Ward-identity
tests.
\end{evidence}

\section{Imaginary roots and evanescent waves}

An imaginary root of an indefinite Kac--Moody algebra is defined by the root
space bilinear form. It is not an imaginary spatial wavenumber. Identifying the
two square roots without a map between the bilinear forms conflates algebraic
classification with a dispersion relation. The repository supplies no such
map, so the proposed \enquote{evanescent threshold} and \enquote{superforce
limit} are unsupported.

\section{Gravitoelectromagnetism}

Gravitoelectromagnetism is a controlled analogy inside general relativity, not
a direct coupling between electromagnetic and gravitational fields. In the
linear perturbation construction, a Minkowski background is perturbed by
$h_{\mu\nu}$, the trace-reversed perturbation obeys a sourced wave equation
after a transverse gauge condition is imposed, and the gravitoelectric and
gravitomagnetic fields are defined from the resulting potentials
\citep{mashhoon2008gravitoelectromagnetism}. The Lorentz-like test-body form is
therefore conditional on the weak-field expansion, the stated gauge and source
definitions, and the retained velocity order.

The analogy does not create a new electromagnetic--gravitational coupling,
does not replace the nonlinear Einstein equations, and does not establish a
propulsion mechanism. A stronger claim requires a covariant action with the
proposed interaction, consistent source and test-body variables, recovery of
the general-relativistic limit, and a controlled prediction that ordinary
electromagnetic coupling, vibration, and thermal drift do not explain.

\section{Materials and zero-point-energy claims}

The repository retains manuscripts concerning tourmaline, vibrations, and
zero-point-energy interpretations. The corpus includes conventional empirical
material data, including tourmaline electrical and thermal measurements, and
standard first-principles calculations of zero-point vibration. This evidence
falsifies an absolute claim that the source collection is wholly nonempirical.
It does not test the proposed enhancement, excess-energy effect, or new
coupling. Those claims require calibrated raw measurements, a complete energy
balance, blinded controls, uncertainty budgets, and independent replication.
The checked-in corpus contains no such discriminating evidence.

The experimental-admission registry converts that boundary into operational
packages. Measurable proposals receive observables, matched controls, complete
energy channels, acceptance criteria, and rejection criteria. Undefined
analogies receive no pretend observable and remain excluded.

\begin{table}[htbp]
  \centering
  \caption{Operational content of the physical-claim admission packages.}
  \label{tab:experimental-admission-packages}
  \small
  \begin{tabularx}{\textwidth}{>{\raggedright\arraybackslash}p{0.43\textwidth}Yrrrr}
    \toprule
    Claim & Outcome & Obs. & Ctrls. & Channels & Missing \\
    \midrule
    Scalar field curvature mechanism & rejected & 3 & 3 & 6 & 4 \\
ZPE energy harvesting & rejected & 3 & 3 & 8 & 4 \\
Time crystal energy amplification & rejected & 3 & 3 & 9 & 4 \\
Tourmaline novel energy effect & bounded & 3 & 3 & 9 & 4 \\
Origami fold merge operator & excluded & 0 & 0 & 0 & 4 \\
Consciousness resonance & excluded & 0 & 0 & 0 & 4 \\
\bottomrule

  \end{tabularx}
\end{table}
